% Generated from locale/tr/content/computability/recursive-functions/pr-functions.tex; do not hand-edit.


\olfileid[tr]{cmp}{rec}{prf}
\olsection{İlkel Özyinelemeli Fonksiyonlar}

İlkel özyineleme ve bileşke kullanarak var olan fonksiyonlardan yeni
fonksiyonları nasıl tanımlayabildiğimizi yeniden kaydedelim.

\begin{defn}
  \ollabel{defn:primitive-recursion}
  $f$ bir $k$-yerli fonksiyon ($k\geq1$), $g$ ise $(k+2)$-yerli bir
  fonksiyon olsun. $f$ ile $g$'den \emph{ilkel özyinelemeyle tanımlanan}
  fonksiyon, aşağıdaki denklemlerle verilen $(k+1)$-yerli $h$ fonksiyonudur:
  \begin{align*}
    h(x_0,\dots,x_{k-1},0) & = f(x_0,\dots,x_{k-1}) \\
    h(x_0,\dots,x_{k-1},y+1) & =
      g(x_0,\dots,x_{k-1},y,h(x_0,\dots,x_{k-1},y)).
  \end{align*}
\end{defn}

\begin{defn}
  \ollabel{defn:composition}
  $f$ bir $k$-yerli fonksiyon; $g_0,\dots,g_{k-1}$ ise her biri $n$-yerli
  olan $k$ fonksiyon olsun. $f$ ile $g_0,\dots,g_{k-1}$ fonksiyonlarından
  \emph{bileşkeyle tanımlanan} fonksiyon, aşağıdaki eşitlikle verilen
  $n$-yerli $h$ fonksiyonudur:
  \[
    h(x_0,\dots,x_{n-1})=
    f(g_0(x_0,\dots,x_{n-1}),\dots,g_{k-1}(x_0,\dots,x_{n-1})).
  \]
\end{defn}

$\Succ$ ile izdüşüm fonksiyonlarının yanı sıra
\[
  \Proj{n}{i}(x_0,\dots,x_{n-1})=x_i,
\]
her doğal sayı $n$ ve $i<n$ için $\Zero(x)=0$ fonksiyonunu da ilkel
özyinelemeli fonksiyonlar arasına katacağız.

\begin{defn}
  İlkel özyinelemeli fonksiyonlar kümesi, $\Nat^n$'den $\Nat$'ye giden ve
  aşağıdaki maddelerle tümevarımsal olarak tanımlanan fonksiyonlar kümesidir:
  \begin{enumerate}
  \item $\Zero$ ilkel özyinelemelidir.
  \item $\Succ$ ilkel özyinelemelidir.
  \item Her izdüşüm fonksiyonu $\Proj{n}{i}$ ilkel özyinelemelidir.
  \item $f$ ilkel özyinelemeli $k$-yerli bir fonksiyon ve
    $g_0,\dots,g_{k-1}$ ilkel özyinelemeli $n$-yerli fonksiyonlarsa,
    $f$'nin $g_0,\dots,g_{k-1}$ ile bileşkesi ilkel özyinelemelidir.
  \item $f$ ilkel özyinelemeli $k$-yerli bir fonksiyon ve $g$ ilkel
    özyinelemeli $(k+2)$-yerli bir fonksiyonsa, $f$ ile $g$'den ilkel
    özyinelemeyle tanımlanan fonksiyon ilkel özyinelemelidir.
  \end{enumerate}
\end{defn}

\begin{explain}
Daha kısa söylersek ilkel özyinelemeli fonksiyonlar kümesi, $\Zero$,
$\Succ$ ve izdüşüm fonksiyonları $\Proj{n}{j}$ öğelerini içeren ve bileşke
ile ilkel özyineleme altında kapalı olan en küçük kümedir.

İlkel özyinelemeli fonksiyonlar kümesi ``aşamalar'' aracılığıyla da
anlatılabilir. $S_0$, başlangıç fonksiyonlarının, yani $\Zero$, $\Succ$ ve
izdüşümlerin kümesi olsun. Bunlar $0$. aşamadaki ilkel özyinelemeli
fonksiyonlardır. Bir $S_i$ aşaması tanımlandıktan sonra $S_{i+1}$, $S_i$
içinde bulunan fonksiyonlara tek bir bileşke veya ilkel özyineleme işlemi
uygulanarak elde edilen bütün fonksiyonların kümesi olsun. O zaman
\[
  S=\bigcup_{i\in\Nat}S_i
\]
bütün ilkel özyinelemeli fonksiyonların kümesidir.
\end{explain}

$\Add$ fonksiyonunun ilkel özyinelemeli olduğunu doğrulayalım.

\begin{prop}
  Toplama fonksiyonu $\Add(x,y)=x+y$ ilkel özyinelemelidir.
\end{prop}

\begin{proof}
$\Add$ için, iki $f$ ve $g$ fonksiyonu cinsinden
\olref{defn:primitive-recursion} biçimine uyan bir ilkel özyinelemeli tanımımız
zaten vardır:
\begin{align*}
  \Add(x_0,0) & = f(x_0)=x_0 \\
  \Add(x_0,y+1) & = g(x_0,y,\Add(x_0,y))=\Succ(\Add(x_0,y)).
\end{align*}
Dolayısıyla $f$ ile $g$ ilkel özyinelemeliyse $\Add$ de öyledir.
$f(x_0)=x_0=\Proj{1}{0}(x_0)$ ve izdüşüm fonksiyonları ilkel özyinelemeli
sayıldığından $f$ ilkel özyinelemelidir. $g$ ise
\[
  g(x_0,y,z)=\Succ(z)
\]
ile tanımlanan üç yerli fonksiyondur. Bu ifade henüz $g$'nin ilkel
özyinelemeli olduğunu göstermez; çünkü $g$ ile $\Succ$ tam olarak aynı
fonksiyon değildir: $\Succ$ bir yerli, $g$ ise üç yerlidir. Fakat $g$'yi
``resmî olarak'' bileşkeyle şöyle tanımlayabiliriz:
\[
  g(x_0,y,z)=\Succ(\Proj{3}{2}(x_0,y,z)).
\]
$\Succ$ ile $\Proj{3}{2}$ ilkel özyinelemeli sayıldığından, ilkel
özyinelemeli fonksiyonlardan bileşkeyle tanımlanan $g$ de ilkel
özyinelemelidir.
\end{proof}

\begin{prop}
  \ollabel{prop:mult-pr}
  Çarpma fonksiyonu $\Mult(x,y)=x\cdot y$ ilkel özyinelemelidir.
\end{prop}

\begin{proof}
  Alıştırma.
\end{proof}

\begin{prob}
  \olref[cmp][rec][prf]{prop:mult-pr} önermesini, $\Mult$ fonksiyonunun
  ilkel özyinelemeli tanımını
  \olref[cmp][rec][prf]{defn:primitive-recursion} içinde gereken biçime
  getirip karşılık gelen $f$ ile $g$ fonksiyonlarının ilkel özyinelemeli
  olduğunu göstererek ispatlayın.
\end{prob}

\begin{ex}
İşte ilkel özyinelemeli tanımın ilk örneği:
\begin{align*}
  h(0) & = 1 \\
  h(y+1) & = 2\cdot h(y).
\end{align*}
$k=0$ olduğundan bu fonksiyon \olref{defn:primitive-recursion} içinde gereken
biçime uymaz. Tanım ayrıca $1$ ve $2$ sabitlerini içerir. İlk sorunu aşmak
için etkisiz bir argüman ekleyip $h'$ fonksiyonunu tanımlayalım:
\begin{align*}
  h'(x_0,0) & = f(x_0)=1 \\
  h'(x_0,y+1) & = g(x_0,y,h'(x_0,y))=2\cdot h'(x_0,y).
\end{align*}
$f(x_0)=1$ fonksiyonu $\Succ$ ile $\Zero$'dan bileşkeyle tanımlanabilir:
$f(x_0)=\Succ(\Zero(x_0))$. $g$ fonksiyonu, $g'(z)=2\cdot z$ ile
izdüşümlerden bileşkeyle tanımlanabilir:
\begin{align*}
  g(x_0,y,z) & = g'(\Proj{3}{2}(x_0,y,z))
  \intertext{$g'$ ise bileşkeyle şöyle tanımlanabilir:}
  g'(z) & = \Mult(g''(z),\Proj{1}{0}(z))
  \intertext{ve}
  g''(z) & = \Succ(f(z)),
\end{align*}
burada $f$ yukarıdaki gibidir: $f(z)=\Succ(\Zero(z))$. Artık $h'$ elimizde
olduğuna göre bileşkeyi bir kez daha kullanarak
$h(y)=h'(\Proj{1}{0}(y),\Proj{1}{0}(y))$ diyebiliriz. Böylece $h$, temel
fonksiyonlardan bir bileşke ve ilkel özyineleme dizisiyle tanımlanır;
dolayısıyla ilkel özyinelemelidir.
\end{ex}

